%%%%%%%%%%%%%%%%%%%%%%
%% consq-no-strings.tex
%%%%%%%%%%%%%%%%%%%%%%

The no-strings rule prohibits a nontrivial charge 
to travel a distance longer than a constant times its size.
We repeat the rule.

\begin{center}
\parbox{.8\textwidth}{
Definition~\ref{defn:no-strings}:
{\em
A string segment is a finitely supported Pauli operator
that creates excitations contained in the union of two finite boxes (anchor regions)
of width $w$.
The string segment is nontrivial 
if the charge contained in one of the boxes is nontrivial.
The distance between the boxes is the length of the string segment.
We say that a model obeys no-strings rule
if the length of any nontrivial string segment of width $w$
is bounded by $\alpha w$ for some constant $\alpha \ge 1$.
}}
\end{center}

\noindent
The cubic code is an extreme case where a point charge cannot move distance $1$
(Theorem~\ref{thm:CubicCode-no-strings}).
It does not mean that isolation or diffusion of charges is impossible;
Theorem~\ref{thm:logarithmic-energy-barrier-for-fractal-operator}
actually gives a local process for the isolation to happen in the translationally invariant case.
The isolation process for a charge has to pay an energy penalty 
upper bounded by the logarithm of the distance from others.

In this chapter, we illustrate the energy landscape of the cubic code Hamiltonian
in a view towards self-correcting quantum memory.
We show that the separation or isolation of charges
requires an energy barrier logarithmically high in the separation distance.
This can be visualized as follows.
Imagine all the energy eigenstates as points on an imaginary ``land,''
and introduce a metric on the land by the minimum number of local operations
one has to apply in order to map one state to another.
Consider a function on the land given by the energy of the energy eigenstate.
Our logarithmic energy barrier means that the energy function, or ``energy landscape,''
has a macroscopic number of local minimums separated by macroscopic energy barriers.
These minimums correspond to low-energy excited states 
in which the separation between defects is approximately the system size.

The energy landscape with the large number of the local minimums 
suggests a possibility of a spin glass phase at a sufficiently low temperature.
Note that there is no quenched disorder in the Hamiltonian~\cite{BinderYoung1986SpinGlass,Weissman1993SpinGlass},
and the glassy feature, if present, would be protected ``topologically''~\cite{BravyiHastingsMichalakis2010stability}.
A spin glass phase can indeed be realized for some classical spin Hamiltonians 
with logarithmic energy barriers such as the model discovered by Newman and Moore~\cite{NewmanMoore1999Glassy}.
Interplay between the topological order and the spin glassiness has been studied recently by
several authors~\cite{Chamon2005Quantum,TsomokosOsborneCastelnovo2010Interplay}.

The logarithmic energy barrier is an optimal bound up to constants
because of Theorem~\ref{thm:logarithmic-energy-barrier-for-fractal-operator}
and \ref{thm:fractal-exists-in-3D}.
An obvious difficulty in proving the lower bound on the energy barrier
is that there are a huge number of paths to isolate a charge.
Relying on scale-invariant nature of the no-strings rule,
we introduce a technique which may be regarded
as a renormalization group in the space of error paths.


Also, we attempt to estimate the length of the separation process:
How many local operations do we need in order to isolate a charge?
Put differently,
how large is a quantum mechanical tunneling probability between two ground states
in presence of local perturbations?
Since the tunneling amplitude would be non-vanishing only for
virtual process that is mediated by an operator acting on the ground space,
the question translates into the theory of error correcting codes
as to find the weight distribution of logical operators.
We do not answer this question at a satisfactory level,
but we find a superlinear lower bound on the code distance assuming the no-strings rule.


\section{Logarithmic energy barrier}
\label{sec:log-barrier}

We consider a regular $D$-dimensional cubic lattice $\Lambda$
with periodic boundary conditions and linear size $L$, that is, $\Lambda=\ZZ_L^D$.
Each site $u\in \Lambda$ is populated by a finite number of qubits.
The class of models we are going to discuss is that of code Hamiltonians
\begin{equation}
H=-\sum_{a=1}^{M} G_a,
\label{eq:stabH}
\end{equation}
where each term $G_a$  is a multi-qubit Pauli operator 
(a tensor product of $I,X,Y,Z$ with an overall $\pm 1$ sign)
and different terms commute with each other;
$G_a G_b =G_b G_a$ and $G_a^2=I$.
We assume that each generator $G_a$ acts nontrivially (by $X,Y$ or $Z$)
only on a set of qubits located at vertices of an elementary cube.
It is allowed to have more than one generator per cube.
Any short-range stabilizer Hamiltonian can be written in this form by performing
a coarse-graining of the lattice.
We continue to assume that $H$ is frustration-free%
\footnote{This is always the case for independent generators $G_a$. 
Since our goal is to obtain a lower bound on the
energy barrier, we can assume that the generators are independent, although it does not
play any role in our analysis.},
The Hamiltonian may or may not be translation-invariant.

Consider any multi-qubit Pauli operator $E$.
A state $\psi = E\,  \psi_0 $
is an excited eigenstate of $H$. Obviously, $G_a \, \psi=\pm \psi$
where the sign depends on whether $G_a$ commutes (plus)
or anticommutes (minus) with $E$.
Any flipped generator ($G_a\, \psi=-\psi$) will be referred to as a {\bf defect}.
Different defects may occupy the same elementary cube.
It should be emphasized that 
a {\em configuration of  defects}, called a {\bf syndrome}, in $\psi$ is the same for all ground states $\psi_0$.
An eigenstate with $m$ defects has energy $2m$ above the ground state.
For brevity, we shall use the term {\bf vacuum} for a  ground state of $H$
whenever its choice  is not important.
A Pauli operator $E$ whose action on the vacuum creates no defects is either a stabilizer ($E\in \mathcal G$),
or a {\bf logical operator} ($E\notin \mathcal G$, but $E$ commutes with $\mathcal G$).
In the former case any ground state of $H$ is invariant under $E$.
In the latter case $E$ maps some ground state of $H$ to an orthogonal ground state.

A Hamiltonian is  said to have topological order if it
has a degenerate ground state and different ground states are locally indistinguishable.
We shall need a slightly stronger version of this condition
that involves properties of both ground and excited states.
These properties depend on a length scale $L_{tqo}$ that must
be bounded as $L_{tqo}\ge  L^\gamma$ for some constant $\gamma>0$.
(For code Hamiltonians, $L_{tqo}$ corresponds to the code distance.)
Most of stabilizer code Hamiltonians with topological order satisfy our conditions
with $L_{tqo} \sim L$.
Our first topological quantum order condition concerns ground states,
which is a rephrasing of Definition~\ref{defn:tqo}:
\begin{defn}
\label{defn:TQO1}
{\bf TQO1} refers to the following condition: 
If a Pauli operator $E$ creates no defects when applied to the vacuum
and its support can be enclosed by a cube of linear size $L_{tqo}$ then $E$
is a stabilizer,  $E\in \mathcal G$.
\end{defn}
Our second TQO condition concerns excited states.
A cluster of defects $S$ will be called {\bf neutral} if it can be created
from the vacuum by a Pauli operator $E$ whose support can be enclosed
by a cube of linear size $L_{tqo}$ without creating any other defects.
Otherwise we say that $S$ is a {\bf charged} cluster.
Given a region $A \subseteq \Lambda$ we shall use a notation
$\mathcal B _r(A)$ for the $r$-neighborhood of $A$, that is, a set of all points that have distance at most $r$ from $A$.
Here and below we use $\ell_{\infty}$-distance on $\ZZ_L^D$.
We shall need the following condition saying that neutral clusters of defects can be
created from the vacuum locally.%
\footnote{If a lattice has a boundary, charged defects
might be created locally on the boundary, as it is the case for the planar version of the toric
code. This is the reason why we restrict ourselves to periodic boundary conditions.}
\begin{defn}
\label{defn:TQO2}
{\bf TQO2} refers to the following condition:
Let $S$ be a neutral cluster of defects 
and $C_{min}(S)$ be the smallest cube that encloses $S$.
Then $S$ can be created from the vacuum 
by a Pauli operator supported on $\mathcal B_1(C_{min}(S))$.
\end{defn}
%
%
\begin{rem}
\label{rem:TransInv-strongTQO}
A translationally invariant exact code Hamiltonian satisfies
both of our topological order conditions.
In particular, the cubic code satisfies the present TQO conditions.
[~$\because$
TQO1 is immediate from Lemma~\ref{lem:local-tqo=exact},
and TQO2 follows from an argument similar to the proof of Lemma~\ref{lem:local-tqo=exact}.
Indeed, given a neutral cluster $e$ of defects,
we have an equation $e = \epsilon p$ for some $p$ of the Pauli module.
$p$ is computed by the standard division algorithm applied to the columns of
the excitation map $\epsilon$.
If $e$ is centered at origin, then the degree of $p$ does not exceed
that of $e$, which implies the second TQO condition.%
]
\end{rem}


Let us consider a process of building a cluster of defects $S$ (syndrome) from the vacuum.
It can be described by an {\bf error path} 
--- a finite sequence of local Pauli errors  $E_1,\ldots,E_T$ such that
$E = E_T \cdots E_2 E_1$ creates $S$ from the vacuum.
For simplicity, we assume that each local error  $E_t$ 
is a single-qubit Pauli operator $\sigma^x$, $\sigma^y$, or $\sigma^z$.
Applying this sequence of errors to a ground state $\psi_0$
generates a sequence of states $\{ \psi(t)\}_{t=0,\ldots,T}$,
where and $\psi(T)=E \, \psi_0$ is the excited state of defect configuration $S$.
We will say that $S$ has {\bf energy barrier $\omega$}
if for any Pauli operator $E$ that creates $S$ from the vacuum 
and for any error path implementing $E$,
at least one of the intermediate states has more than $\omega$ defects.
Note that we do not impose any restriction on the length of the path $T$ as long as it is finite.
In particular, one and the same error may be repeated in the error path several times
at different time steps.
Similarly, we consider the energy barrier for a logical operator $\overline{P}$;
we say that a logical operator $\overline{P}$ has {\bf energy barrier $\omega$}
if for any error path implementing $\psi_0 \mapsto \overline P \, \psi_0$
at least one of the intermediate states $\psi(t)$ has more than $\omega$ defects.


For any integer $p\ge 0$, define a {\bf level-$p$ unit of length $\xi(p)$} as%
\footnote{The choice of the constant $10$ is somewhat arbitrary.
We do not try to optimize constants in our proof.}
%%%%%%%%%%%%%%%%%%%%%%%%%%%%%%%%%%%%%%%%%%%%%%%%
\begin{equation}
\xi(p)=(10\alpha)^p, \quad p=0,1,\ldots.
\end{equation}
%%%%%%%%%%%%%%%%%%%%%%%%%%%%%%%%%%%%%%%%%%%%%%%%
Let $S$ be any non-empty syndrome.
Recall that each defect in $S$ can be associated with some elementary cube of the lattice.
\begin{defn}
\label{defn:psparse}
A syndrome $S$ is said to be {\bf sparse at level $p$} 
if the set of elementary cubes occupied by the defects in $S$
can be partitioned into a disjoint union of clusters
such that each cluster has diameter at most $\xi(p)$
and any pair of distinct clusters combined together has diameter larger
than $\xi(p+1)$. Otherwise, $S$ is {\bf non-sparse at level $p$}.
\end{defn}
\noindent
For example, suppose all defects in $S$ occupy the same elementary cube.
Since an elementary cube has diameter $1$, such a syndrome $S(t)$ is sparse at any level $p \ge 0$.
If $S$ occupies a pair of adjacent cubes, $S(t)$ is sparse at any level $p\ge 1$,
but is non-sparse at level $p=0$.
Note that the partition of $S$ into clusters required for level-$p$ sparsity
is unique whenever it exists.
The non-sparsity provides a lower bound on the number of defects in a cluster
as follows.
\begin{lem}
\label{lemma:counting}
A non-empty syndrome $S$ that is non-sparse at all levels $q=0,\ldots,p$
contains at least $p+2$ defects.
\end{lem}
\begin{proof}
Let $C^{(0)}_1,\ldots, C^{(0)}_g$ be elementary cubes occupied by $S$.
Obviously, $S$ contains at least $g$ defects.
Since $S$ is non-empty and non-sparse at level $0$,
we have $g \ge 2$ and there exists
a pair of cubes $C_a^{(0)}, C_b^{(0)}$ such that the union
$C_a^{(0)} \cup C_b^{(0)}$ has diameter at most $\xi(1)$.
Combining the pair $C_a^{(0)}, C_b^{(0)}$  into a single cluster,
we obtain a partition $S = C_1^{(1)} \cup \cdots \cup C_{g-1}^{(1)}$
where each cluster $C_a^{(1)}$ has diameter at most $\xi(1)$.
Since $S$ is non-sparse at level $1$, we have $g-1 \ge 2$,
and there exists a pair of clusters $C_a^{(1)}, C_b^{(1)}$ such that the union
$C_a^{(1)}\cup C_b^{(1)}$ has diameter at most $\xi(2)$.
Combining the pair $C_a^{(1)}, C_b^{(1)}$ into a single cluster and
proceeding in the same way we arrive at $g \ge p+2$.
\end{proof}

A configuration of defects created by applying a Pauli operator $E$ to the vacuum
will be called a {\em syndrome caused by $E$}.
The process of building up a logical operator $\overline{P}$ 
by a sequence of local errors $E_1,\ldots,E_T$ can be described by
a {\bf syndrome history} $\{S(t)\}_{t=0,\ldots,T}$.
Here $S(t)$ is the syndrome caused by the product $E_t \cdots E_2 E_1$,
a partial implementation of $\overline{P}$ up to a step $t$.
The following concerns the size of local errors
\begin{lem}
\label{lem:large-support-for-nontrivial-transition}
Let $Q_j$ be Pauli operators causing  a chain of transitions
\[
\mathrm{vac} \xrightarrow{Q_1} S_1 \xrightarrow{Q_2}  S_2 \xrightarrow{Q_3} \ldots \xrightarrow{Q_r} S_r
\xrightarrow{Q_{r+1}} \mathrm{vac}.
\]
Let $P_j$ be some Pauli operator creating the syndrome $S_j$ from the vacuum.
Suppose the support of any operator $P_j$ and any operator $Q_j$
can be enclosed by $n$ or less cubes of linear size $R$ such that $4nR < \ltqo$.
Then, the product $\overline{Q} = Q_1\cdots Q_r Q_{r+1}$ is a stabilizer.
\end{lem}
\begin{proof}
Let $\psi_0$ be any ground state.
Define a sequence of states
\begin{align*}
\psi(1) &= P_1 Q_1\cdot \psi_0, \\
\psi(j+1) &=  (P_j P_{j+1}) Q_{j+1} \cdot \psi(j) \quad \text{ for $j=1,\ldots,r-1$},\\
\psi(r+1) &= Q_{r+1} P_{r} \cdot \psi(r).
\end{align*}
Obviously,
\begin{align*}
\psi(j) &= \pm P_j \cdot (Q_1 \cdots Q_j) \cdot  \psi_0 \quad \text{for $j=1,\ldots,r$} \\
\psi(r+1) &= \pm \overline{Q}\cdot \psi_0.
\end{align*}
It follows that all states $\psi(j)$ are ground states,
and the transition from $\psi(j)$ to $\psi(j+1)$ 
can be caused by a Pauli operator
\[
O_j = P_j P_{j+1} Q_{j+1}.
\]
Let $M_j$ be the support of $O_j$. 
By assumption, $M_j$ can be enclosed by at most $3n$ cubes of linear size $R$.
If $M_j$ is a connected set,
i.e., one can connect any pair of qubits from $M_j$ by a path $(u_1,\ldots,u_l)$
such that the distance between $u_a$ and $u_{a+1}$ is $1$,
then $M_j$ can be enclosed by a single cube of linear size at most $3nR$.
That $3nR < \ltqo$ implies $O_j$ is a stabilizer by the topological order condition.
Generally, $M_j$ consists of several disconnected components $M_j^\alpha$,
such that the distance between any pair of distinct components is at least $2$.
The restriction of $O_j$ on a connected component 
commutes with any stabilizer generator,
and is supported in a box of linear size $3nR$,
which is smaller than $\ltqo$ by assumption.
Therefore, the restrictions are stabilizers,
and their product $O_j$ is also a stabilizer.
In other words, $\psi(j+1)=\pm \psi(j)$ for all $j$.
It means that $\overline{Q}\, \psi_0 = \pm \psi_0$ for any ground state $\psi_0$.
We conclude that $\overline{Q}$ is a stabilizer.
\end{proof}


The no-strings rule says that an isolated charged defect belonging to some sparse syndrome
cannot be moved further than distance $\alpha$ away by a sequence of local errors.
Since the no-strings rule is scale invariant,
it may be applied to a coarse-grained lattice to show that isolated
charged clusters cannot be moved further than distance $\alpha \xi(p)$ away.
In order to exploit the scale invariance,
we define a {\bf level-$p$ syndrome history}
as a subsequence of the original syndrome history $\{S(t)\}_{t=0,\ldots,T}$
that includes only those syndromes $S(t)$ that are non-sparse
at all levels $q=0,\ldots,p-1$.
The level-$0$ syndrome history includes all syndromes $S(t)$,
the level-$1$ syndrome history omits $S(t)$ that is sparse at level 0,
and so on.
When $S(t')$ and $S(t'')$ are a consecutive pair of level-$p$ syndromes,
we define a {\bf level-$p$ error $E$} connecting $S(t')$ and $S(t'')$
as the product of all single-qubit errors $E_j$ that occurred between $S(t')$ and $S(t'')$.
Level-$p$ errors are represented by horizontal arrows on Fig.~\ref{fig:RG}.
Note that we do not have any bound on the number
of single-qubit errors $E_j$ in the interval between $S(t')$ and $S(t'')$.
In the worst case, $E$ could act nontrivially on every qubit in the system.
%
A main technical lemma of this chapter below asserts, loosely speaking, that
if a syndrome history does not have a deep enough non-sparse hierarchy,
any error path is equivalent to one that is localized around the defects.
% We would like to show that $E$ can be regarded as an approximately local error
% on a coarse-grained lattice characterized by the unit of length $\xi(p)$.


\begin{lem}
\label{lemma:RG2}
Let $S'= S(t')$ and $S''= S(t'')$ be a consecutive pair of syndromes in the level-$p$ syndrome history
of a Hamiltonian obeying the no-strings rule.
Let $E$ be the product of all errors $E_j$ that occurred between $S'$ and $S''$.
Suppose that any $S(t)$ contains at most $m$ defects.
If 
\[
16 m \xi(p) < L_{tqo},
\]
then there exists an error $\tilde{E}$  supported on $\mathcal B _{\xi(p)}(S'\cup S'')$
such that $E\tilde{E}$ is a stabilizer.
\end{lem}

\begin{figure}[tb]
\centerline{\includegraphics[height=4cm]{rg_scheme5}}
\caption{Renormalization group technique 
used to prove a logarithmic  lower bound on the energy barrier for logical operators.
Horizontal axis represents time. 
Vertical axis represents RG level $p=0,1,\ldots,p_{max}$.
A sequence of level-$0$ errors (single-qubit Pauli operators) 
implementing a logical operator $\overline{P}$
defines a level-$0$ syndrome history (yellow circles)
that consists of sparse (S) and non-sparse (D) syndromes.
The history begins and ends with the vacuum ($0$).
For any level $p\ge 1$ we define a level-$p$ syndrome history
by retaining only non-sparse syndromes at the lower level.
A syndrome is called non-sparse at level $p$ if it
cannot be partitioned into clusters of size $\le (10\alpha)^p$
separated by distance $\ge (10\alpha)^{p+1}$,
where $\alpha$ is a constant coefficient from the no-strings rule.
Each  level-$p$ error (horizontal arrows) connecting syndromes $S',S''$
is equivalent to the product of all level-$(p-1)$ errors between $S',S''$
modulo a stabilizer.
We prove that these stabilizers can be chosen such that
level-$p$ errors act on $2^{O(p)}$ qubits.
Since at the highest  level $p=p_{max}$
a single level-$p$ error is a logical operator,
one must have $p_{max}=\Omega(\log{L})$.
We prove that level-$p$ non-sparse syndromes contain $\Omega(p)$ defects
which implies that at least one syndrome at level $p=p_{max}-2$
consists of $\Omega(\log{L})$ defects.}
\label{fig:RG}
\end{figure}

\begin{proof}
The proof is by induction in $p$.
When $p=0$, $E=E_j$ is a single-qubit error.
If the qubit acted on by $E$ does
not belong to $\mathcal B_{1}(S'\cup S'')$, one must have $S'=S''$.
It means that $E$ is a single-qubit error with a trivial syndrome.
The topological order condition implies that $E$ is a stabilizer.
Choosing $\tilde{E}=I$ proves the lemma for $p=0$.

Suppose the assertion is true for some level $p \ge 0$.
Let $S'=S(t')$ and $S''=S(t'')$ be consecutive syndromes in the level-$(p+1)$
history. Consider first the trivial case when
$S'=S(t')$ and $S''=S(t'')$ are also consecutive syndromes in the level-$p$ history.
Then $S'$ and $S''$ are connected by a single%
\footnote{%
The word `single' does not necessarily mean a single blob of errors supported in a small ball; 
it just means a single arrow in the level-$p$ history.
For instance, suppose that a model permits neutral point defects.
Two neutral defects separated by distance 1 is a non-sparse at level 0 and 1.
If we annihilate one of them, put sparsely many neutral defects,
and finally put a neutral defect adjacent to one of the defects,
then the whole process is described by a consecutive pair in the level 1 history.
The `single' level-1 error in this case consists of many components.%
}
level-$p$ error $E$ which, by induction hypothesis, has support on $\mathcal B _{\xi(p)}(S'\cup S'')$
modulo stabilizers.
The latter is contained in $\mathcal B_{\xi(p+1)}(S'\cup S'')$ which proves the induction step.

The nontrivial case is when there is at least one level-$p$ syndrome between $S'$ and $S''$.
The interval of the level-$p$ syndrome history between $S'$ and $S''$ can be represented
(after properly redefining the time variable $t$) as
\[
S'\xrightarrow{E_\text{lead}}
S(1) \xrightarrow{E_1}
S(2) \xrightarrow{E_2}
\cdots
 \xrightarrow{E_{\tau-1}}
S(\tau) \xrightarrow{E_\text{tail}}
S''.
\]
Here all syndromes $S(1),\ldots,S(\tau)$ are sparse at the level $p$ and
all transitions are caused by level-$p$ errors.
The sparsity implies that
the set of elementary cubes occupied by  $S(t)$ has a unique
partition into a disjoint union of clusters $C_a(t)$ such that
each cluster has diameter at most $\xi(p)$ and the distance
between any pair, if any, of clusters is at least
\begin{align*}
\mathrm{dist}(C_a(t),C_b(t)) & \ge   \xi(p+1)-2\xi(p) \ge (10\alpha -2)\xi(p)  \ge   8\alpha \xi(p).
\end{align*}
Represent any intermediate syndrome as a disjoint union
\begin{equation}
\label{S1+}
S(t)=S^c(t) \cup S^n(t), \quad t=1,\ldots,\tau,
\end{equation}
where $S^c(t)$ and $S^n(t)$ include all charged and all neutral clusters $C_a(t)$, respectively.
Let $g$ be  the number of clusters in $S^c(t)$. We claim  that $g$  does not depend on $t$.
Indeed, since a level-$p$ error $E_t$ acts on $\xi(p)$-neighborhood of $S(t) \cup S(t+1)$
by the induction hypothesis,
the sparsity condition implies that $E_t$
cannot create or annihilate a charged cluster $C_a(t)$ from the vacuum, 
or map a charged cluster  to a neutral cluster and vice versa.
The same argument shows that each  cluster $C_a(t)\subseteq  S^c(t)$ can `move' 
at most by $\xi(p)$ per time step,
that is, we can parameterize
\[
S^c(t)=C_1(t) \cup \ldots \cup C_g(t)
\]
such that
a `world-line' of the $a$-th charged cluster
obeys the continuity condition
\begin{equation}
\label{cont+}
\mathrm{dist}(C_a(t+1),C_a(t)) \le \xi(p).
\end{equation}
We can now use the no-strings rule to show that
all charged clusters are `locked' near their
initial positions, so that their world-lines
are essentially parallel to the time axis.  More precisely, we claim that
\begin{equation}
\label{locking+}
\mathrm{dist}(C_a(t), C_a(1)) \le \alpha  \xi(p)  \quad \text{for all $1\le t \le \tau$}.
\end{equation}
Indeed, suppose Eq.~\eqref{locking+} is false for some $a$.
Using the continuity Eq.~\eqref{cont+},
one can find a time step $t_1$
such that $\mathrm{dist}(C_a(t_1), C_a(1))>\alpha \xi(p)$ and
$\mathrm{dist}(C_a(t), C_a(1)) \le \alpha \xi(p)$ for all $1\le t<t_1$.
Let $E_{close}$ be the product of all level-$p$ errors $E_j$
that occurred between $S(1)$ and $S(t_1)$ within distance $(2+\alpha)\xi(p)$
from $C_a(1)$. Since all intermediate syndromes are sparse at level $p$,
the net effect of $E_{close}$ is to annihilate the  charged cluster  $C_a(1)$
and create the charged cluster $C_a(t_1)$.
Equivalently, applying $E_{close}$ to the vacuum creates a pair of charged clusters
$C_a(1)$ and $C_a(t_1)$.
However, this contradicts to the no-strings rule since
$C_a(1)$ and $C_a(t_1)$ have linear size at most $\xi(p)$ while
the distance between them is greater than $\alpha \xi(p)$.
Thus we have proved Eq.~\eqref{locking+}.

We say that $\vec{x}\in \Lambda$ is {\em close to $S'$} if $\vec{x}\in \mathcal B_{\xi(p+1)}(S')$,
and $\vec{x}\in \Lambda$ is {\em close to $S''$} if $\vec{x}\in \mathcal B_{\xi(p+1)}(S'')$.
%
Let $E_t$ be the level-$p$ error causing the transition from  $S(t)$ to $S(t+1)$,
where $t=1,\ldots,\tau-1$.
By induction hypothesis, we may assume that the support of $E_t$ is in $\mathcal B_{\xi(p)}(S(t) \cup S(t+1))$.
Let $E^c_t$ be the restriction of $E_t$ onto $\mathcal B_{\xi(p)}(S^c(t)\cup S^c(t+1))$,
and $E^n_t$ be the restriction of $E_t$ onto $\mathcal B_{\xi(p)}(S^n(t)\cup S^n(t+1))$.
The sparsity of $S(t)$ implies that
\begin{equation}
\label{EcEn+}
E_t = E^c_t \cdot E^n_t.
\end{equation}
We claim that {\em any error $E^c_t$ is close to $S'$}.
Indeed, each cluster in $S^c(1)$ is within distance $2\xi(p)$ from $S'$ 
since, otherwise, a single level-$p$ error $E_\text{lead}$,
supported in $\mathcal B_{\xi(p)}(S' \cup S(1) )$ by induction hypothesis (modulo stabilizers),
would be able to create a charged cluster from the vacuum.
Using Eq.~\eqref{locking+}, we infer that 
$C_a(t) \subseteq \mathcal B_{(2+\alpha)\xi(p)}(S')$ for all $a=1,\ldots,g$.
Therefore, $E^c_t$ is close to $S'$.


We wish to find a ``localized'' leading error $\tilde{E}_\text{lead}$
that maps the syndrome $S'$ to $S^c(1)$ such that the support of
$\tilde{E}_\text{lead}$ is close to $S'$.
For each neutral cluster $C \in S^n(1)$ of diameter at most $\xi(p)$,
let $O'(C)$ be a Pauli operator creating $C$ from the vacuum.
Because of our TQO2, we can choose $O'(C)$ to be supported in $\mathcal B_1(C)$.
Set
\[
\tilde{E}_\text{lead} = E_\text{lead} \prod_{ C \in S^n(1) } O'(C),
\]
We have seen that $S^c(1)$ is within $(2+\alpha)\xi(p)$-neighborhood of $S'$.
If $E_\text{lead}$ has a disconnected component $E(C)$,
isolated by distance $\xi(p)$ and centered at a neutral cluster $C$ of $S^n(1)$,
then $E(C) O'(C)$ is a stabilizer.
Hence, $\tilde E_\text{lead}$ is close to $S'$ modulo stabilizers.
If $E_\text{lead}$ does not have such a component, $\tilde E_\text{lead}$ is already close to $S'$.
Certainly, $\tilde E_\text{lead}$ maps $S'$ to $S^c(1)$.
We can apply the same rules to the error $E_\text{tail}$ and the syndrome $S^n(\tau)$.
We find a localized error
\[
\tilde{E}_\text{tail} = E_\text{tail} \cdot \prod_{C \in S^n(\tau) } O''(C)
\]
modulo stabilizers
such that $\tilde{E}_\text{tail}$ maps $S^c(\tau)$ to $S''$ 
and the support of $\tilde{E}_\text{tail}$ is close to $S''$.
The operator $O''(C)$ above creates a neutral cluster $C \in S^n(\tau)$ from the vacuum.

We can now define a localized level-$(p+1)$ error $\tilde{E}$
whose support is close to $S'\cup S''$ as
\begin{equation}
\label{tildeE_cluster}
\tilde{E}=\tilde{E}_\text{tail} \cdot E^c_{\tau-1} \cdots E^c_{1} \cdot \tilde{E}_\text{lead}.
\end{equation}
By construction, it describes an error path
\[
S'\xrightarrow{\tilde{E}_\text{lead}}
S^c(1) \xrightarrow{E^c_1}
S^c(2) \xrightarrow{E^c_2}
\cdots
 \xrightarrow{E^c_{\tau-1}}
S^c(\tau) \xrightarrow{\tilde{E}_\text{tail}}
S''.
\]
It remains to check that $E \cdot \tilde{E}$ is a stabilizer.
Combining Eq.~\eqref{EcEn+} and
Eq.~\eqref{tildeE_cluster} we
conclude that
\[
E\cdot \tilde{E} = \pm (\tilde{E}_\text{lead} E_\text{lead}) \cdot E^n_1 \cdots E^n_{\tau-1} \cdot (\tilde{E}_\text{tail} E_\text{tail}).
\]
Applying  $E\cdot \tilde{E}$ to the vacuum generates the following chain of transitions:
\begin{align}
\label{chain1_cluster}
\mathrm{vac} \xrightarrow{\tilde{E}_\text{lead} E_\text{lead}}
S^n(1) \xrightarrow{E^n_1}
S^n(2) \xrightarrow{E^n_2} \cdots 
 \xrightarrow{E^n_{\tau-1}}
S^n(\tau) \xrightarrow{\tilde{E}_\text{tail} E_\text{tail}}
\mathrm{vac}
\end{align}
Each syndrome $S^n(t)$ consists of at most $m$ neutral clusters of diameter $\xi(p)$, 
i.e., it can be created from the vacuum by an error
whose support can be enclosed by at most $m$ cubes of linear size $2+\xi(p)$,
due to our TQO2.
The first transition $\tilde{E}_\text{lead} E_\text{lead}$
or the last transition $\tilde{E}_\text{tail} E_\text{tail}$
is caused by errors whose support can be enclosed by at most $m$ cubes of linear size $2+\xi(p)$.
All intermediate ones are supported on $\mathcal B_{\xi(p)}(S^n(t)\cup S^n(t+1))$.
Hence their support can be enclosed by at most $2m$ cubes of linear size $2+\xi(p)$.
Now the statement that $E\cdot \tilde{E}$ is a stabilizer follows from Lemma~\ref{lem:large-support-for-nontrivial-transition}.
\end{proof}

Now we state two theorems that apply to any stabilizer Hamiltonian Eq.~\eqref{eq:stabH}
on a $D$-dimensional lattice that obeys the topological order condition (TQO1,2)
and the no-strings rule.
The numerical constants for the Hamiltonian are
$\alpha \ge 1$ in the no-strings rule
and $1 \ge \gamma>0$ in the bound $\ltqo \ge L^\gamma$
where $L$ is the linear system size.

\begin{theorem}
\label{thm:lop-log-energy-barrier}
The energy barrier for any logical operator is at least $c \log{L}$ for some constant $c=c(\alpha,\gamma)$.
\end{theorem}
\begin{proof}
Consider a syndrome history of an implementation of a logical operator $E$.
We keep the initial and the final syndromes (the empty syndromes) at all levels;
the syndrome history starts and ends with the vacuum at any level $p$.
It suffices to treat the case where
all intermediate syndromes $S(t)$ are non-empty.
Let $p_{max}$ be the highest RG level, that is, the smallest integer $p \ge 0$
such that a single level-$p$ error $E$ maps the vacuum to itself, see Fig.~\ref{fig:RG}.
Let $m$ be the maximum number of defects in the syndrome history at any given moment.

Suppose that $16m \xi(p_{max}) < \ltqo$.
Then Lemma~\ref{lemma:RG2} applied to the level-$p_{max}$ syndrome history
with $S'=S''=\emptyset$ (vacuum),
would imply $\tilde{E}=I$.
Since $E$ is not a stabilizer by assumption, we must have $16 m \xi(p_{max}) \ge \ltqo$.
Since the syndrome history must contain a syndrome $S(t)$ non-sparse at all levels $0,\ldots,p_{max}-2$,
Lemma~\ref{lemma:counting} implies $m \ge p_{max}$.
And, TQO1 requires $\ltqo \ge L^\gamma$.
Therefore, $ 16(10\alpha)^{2m} \ge 16m(10\alpha)^m \ge 16 m (10\alpha)^{p_{max}} \ge L^\gamma $,
and $m = \Omega(L)$.
\end{proof}

\begin{theorem}
\label{thm:cluster-log-energy-barrier}
Let $S$ be a neutral cluster of defects
containing a charged cluster $S' \subseteq S$ of diameter $r$
such that there are no other defects within distance $R$ from $S'$.
If $r+R < L_{tqo}$,
then the energy barrier for creating $S$ from the vacuum is at least $c \log R$
for some constant $c=c(\alpha)$.
\end{theorem}
\begin{proof}
Let $S$ be a neutral cluster of defects and $E$ be a Pauli operator creating $S$ from the vacuum,
with $S' \subset S$ of diameter $r$ being charged.
Consider a hierarchy of syndrome histories similar to the one shown on Fig.~\ref{fig:RG},
where we now maintain the initial syndrome $\emptyset$ and the final syndrome $S$ for all levels.
Let $p_{max}$ be the highest RG level.
Then a single level-$p_{max}$ error $E$ creates $S$ from the vacuum.
Since there must be a syndrome that is non-sparse for all levels $0,1,\ldots,p_{max}-2$,
Lemma~\ref{lemma:counting} implies $m \ge p_{max}$
where $m$ is the maximum number of defects in the syndrome history.

Suppose $16 m \xi(p_{max}) < \ltqo$ 
Lemma~\ref{lemma:RG2} implies that 
$E$ is the equivalent modulo stabilizers to $\tilde{E}$ 
supported on $\mathcal B _{\xi(p_{max})}(S)$.
If $\xi(p_{max}) < R/4$,
then $\tilde{E}$ must act on two separated regions, one near $S'$ and another far from $S'$.
This means that $S'$ alone can be created by a Pauli operator whose support is enclosed by a cube of linear size $r+R/2$.
Since $r+R < L_{tqo}$, it is contradictory to the assumption that $S'$ is charged.
Therefore, $\xi(p_{max}) \ge R/4$, and $m = \Omega(R)$.
%
If $16 m \xi(p_{max}) \ge \ltqo$, then, since $\ltqo > R$, we also have $m = \Omega(R)$.
\end{proof}



\section{Superlinear code distance}
\label{sec:superlinear-distance}

We have shown that in a process of isolating a charged cluster, there is a logarithmic energy barrier.
The following theorem quantifies how long the process must be.
%
The proof again makes use of renormalization group,
and shows that there is a subset of `fractal dimension' $\gamma > 1$ in the support of $E$,
the operator that makes two separated clusters from the vacuum of which one is charged.
We assume that $E$ has weight minimum possible.

Let $w$ be an odd positive number.
We say a set of sites $C \subset \Lambda$ is a {\bf level-$p$ chunk} if $\mathrm{diam}(C) < w^p$.
A \emph{path} in the lattice is a finite sequence of sites $(u_1,u_2,\ldots,u_n)$ such that $d(u_i,u_{i+1})=1$.
(Recall that we use the $l_\infty$ metric $d$.)
Using paths, we can say whether a set is connected.

\begin{defn}
A connected level-$p$ chunk $C \subseteq S$ is \emph{maximal} with respect to a set of sites $S$
if there exist a connected subset $C^\circ \subseteq C$
and a path $\zeta = (u_1,\ldots,m,\ldots,u_n) \subseteq C^\circ$ satisfying
\begin{enumerate}
 \item[(i)]   $d(u_1,u_n)=w^p-w^{p-1}$,
 \item[(ii)]  $d(u_1,m),d(u_n,m)\ge \frac{w^p-w^{p-1}}{2}$,
 \item[(iii)]  $C^\circ$ contains the connected component of $m$ in $\mathcal{B}_{\frac{w^{p}-w^{p-1}}{2}}(m) \cap S$, and
 \item[(iv)] $C$ contains the connected component of $C^\circ$ in $\mathcal{B}_{\frac{w^{p-1}}{2}}(C^\circ)\cap S$.
\end{enumerate}
\end{defn}

The last two conditions restricts the position of $\zeta$ in $C$ such that $\zeta$ lies sufficiently far from the boundary of $C$.
The site $m$ will be referred to as a {\bf midpoint} of $C$.
Let $S$ be the support of the Pauli operator $E$,
any restriction of which obeys the no-strings rule.

\begin{lem}
Given a path $\zeta$ in $S$ joining $u_1$ and $u_n$ such that $d(u_1,u_n)=lw^p-1$,
there are $l$ disjoint maximal chunks of level $p$ whose midpoints are on $\zeta$.
\label{lem:max_chunks_on_path}
\end{lem}
\begin{proof}
For convenience, we assume that the $z$-coordinates of $u_1$ and $u_n$ are
$0$ and $lw^p-1$, respectively.
Consider $l+1$ planes $P_i$ perpendicular to the $z$-axis, whose $z$-coordinates are $iw^p$ for $i=0,1,\ldots,l$.
In each region between the two consecutive planes $P_{i-1}$ and $P_i$,
there is a subpath $\zeta_i=(u_{j_{i-1}},\ldots,u_{j_{i}})$ such that $d(u_{j_{i-1}},u_{j_{i}})=w^p-1$.
Choose $m_i \in \zeta_i$ such that $d(u_{j_{i-1}},m_i),d(m_i,u_{j_{i}}) \ge \frac{w^p -1}{2}$.
Let $C_i$ be the connected component of $m_i$ within $S \cap \mathcal B_{\frac{w^p}{2}}(m_i)$.
Add, if necessary, some points of $\zeta_i$ to $C_i$ to get a maximally connected $C'_i$.
This $C'_i$ is a maximal chunk of sites with midpoint being $m_i$.
Any two $C'_i$'s are disjoint since each of them lies in a unique region enclosed by $P_{i-1}$ and $P_i$.
\end{proof}


\begin{lem}
For sufficiently large $w$,
a maximal level $(p+1)$ chunk $C$ with respect to $S$
admits a decomposition into $w+1$ or more maximal chunks of level $p$ with respect to $S$.
\label{lem:many_max_in_max}
\end{lem}
\begin{proof}
Recall that $S$ is the support of the Pauli operator $E$, any restriction of which obeys no-strings rule.
Define the \emph{boundary} of a subset $U$ of $S$ to be $\partial U = \mathcal B_1(U) \cap U^c \cap S$.
Then, any subset $U$ of sites with boundary enclosed in a two disjoint regions can be regarded as a string segment.

By the definition of the maximal chunk,
there exists a path $(u_1,\ldots,m,\ldots,u_n)$ in $C^\circ \subseteq C$ such that $d(u_1,u_n)=w^{p+1}-w^p$.
We assume that the $z$-coordinates of $u_1, u_n$ differ by $w^{p+1}-w^p$.
We will show that there are sufficiently long and separated paths in $C^\circ$,
to which we apply Lemma~\ref{lem:max_chunks_on_path} to find $w+1$ maximal chunks of level $p$.
They will lie in $\mathcal B_\frac{w^p}{2}(C^\circ)$, and hence in $C$.

Let $M$ ($N$) be the subset of $S$ consisted of sites
whose $z$-coordinates differ from that of $u_1$ ($u_n$) by at most $\eta w^p$.
First, suppose $\partial C^\circ$ is not contained in $M \cup N$.
Since $u_1 \in M$ and $u_n \in N$, there is a site $s \in C^\circ$
adjacent (of distance 1) to $\partial C^\circ$ such that $d(s,u_1),d(s,u_n) > \eta w^p$.
Furthermore, $d(s,m) \ge \frac{w^{p+1}-w^p}{2}-1$;
otherwise, $C^\circ$ contains a site in the boundary, which is a contradiction.

Consider the shortest network $\mathcal{N}$ of paths in $C^\circ$ connecting four sites $u_1,m,u_n,s$.
(The length of a network of paths is the number of sites in the union of the paths.)
Let $\zeta$ be the shortest path in $\mathcal{N}$ from $u_1$ to $u_n$.
If $s$ is not contained in $\mathcal B_{3w^p}(\zeta)$,
then $\zeta' \subseteq \mathcal{N}$ joining $s$ to a site on $\zeta$
has a subpath $\zeta'' \subseteq \zeta'$ of diameter at least $2w^p$
such that $\zeta''$ is separated from $\zeta$ by $w^p$.
Applying Lemma~\ref{lem:max_chunks_on_path} to $\zeta$ and $\zeta''$,
we find at least $w+1$ maximal chunk of level $p$.
If $m$ is not contained in $\mathcal B_{3w^p}(\zeta)$,
a similar argument reveals at least $w+1$ maximal chunk of level $p$.

Suppose both $s$ and $m$ are contained in $\mathcal B_{3w^p}(\zeta)$.
Observe that $\zeta \setminus ( \mathcal B_{4w^p}(s) \cup \mathcal B_{4w^p}(m) )$
consists of three connected components $\zeta_1,\zeta_2,\zeta_3$,
two of which have diameter $\ge \frac{w^{p+1}-w^p}{2}-8w^p$ and the other has diameter $\ge (\eta - 4)w^p$.
Two distinct $\mathcal B_{\frac{w^p}{2}}(\zeta_i)$ and $\mathcal B_{\frac{w^p}{2}}(\zeta_j)$ ($i,j=1,2,3$)
do not overlap because of the minimality of $\zeta$. Applying Lemma~\ref{lem:max_chunks_on_path},
we find $w + \eta - 21$ maximal chunks of level $p$.
Choosing $\eta > 21$, we get the desired result.

Next, suppose $\partial C^\circ$ is contained in $M \cup N$.
Let $s_M, s_N \in C^\circ \setminus (M\cup N)$ be sites adjacent to $M$ and $N$, respectively.
The separation between $M$ and $N$ is $(w-1-2\eta)w^p$.
If it is greater than $\eta' \alpha w^p$, there must be a $z$-plane $P$ that contains $s_M$ or $s_N$
such that $P \cap C^\circ$ has diameter $> \eta' w^p$; Otherwise, the no-strings rule is violated.
Let $v_1, v_2 \in P \cap C^\circ$ be sites separated by $\eta' w^p$.
The four sites, $u_1,u_n,v_1,v_2$ are sufficiently separated and connected by some paths in $C^\circ$.
Arguing as before, we find $(w -1) + \eta'-16$ maximal chunks of level $p$.
The choice of $\eta' > 17$ and $w > 1+2\eta+\eta'\alpha$ proves the lemma.
\end{proof}


\begin{theorem}
Let $E$ be a Pauli operator creating $S$, a neutral cluster of defects
containing a charged cluster $S' \subseteq S$ of diameter $r$
such that there are no other defects within distance $R$ from $S'$.
If $r+2R < L_{tqo}$, then the weight of $E$ must be $\ge c R^\gamma$ for some constant $\gamma > 1$ and $c$.
\label{thm:separation_charged_defects}
\end{theorem}
\begin{proof}
The support of the minimal Pauli operator $E$ in Theorem~\ref{thm:separation_charged_defects}
must admit a path connecting $S'$ and $S \setminus S'$.
Otherwise, $S'$ can be regarded as being created locally,
and our topological order condition demands the cluster be neutral.
Since the path has length $\ge R$,
Lemma~\ref{lem:max_chunks_on_path} says we have a maximal chunk of level $p$
where $p$ is such that $w^p \le R < w^{p+1}$.
Lemma~\ref{lem:many_max_in_max} implies any maximal chunk of level $p$ must contain
at least $(w+1)^p$ sites.
This concludes the proof with $\gamma=\frac{\log(w+1)}{\log w}>1$.
\end{proof}
A similar argument proves the lower bound $d=\Omega(L^\gamma)$ on the code distance $d$ of the cubic code
since the minimal logical operator must contain a path of length $L$.